\section{Where Values Come From: Prestige, CREDs, and the Cultural-Evolution
Substrate}\label{sec:cultural-evolution}

Humans are, in a fundamental sense, obligate cultural learners: the overwhelming
share of what any individual knows, wants, and treats as worth attending to was
acquired from other people rather than rediscovered from scratch
\citep{henrichGilWhite2001Prestige}.
This dependence runs deeper than the transmission of practical skills or factual
beliefs.
What counts as a \emph{value} — which outcomes are desirable, which actions are
admirable, which objects are sacred or disgusting — is itself downstream of what
gets attended to, imitated, and rehearsed in a community's daily life.
A child who grows up watching elders fast, or observing neighbours donate
conspicuously to shared projects, or noticing which adults attract the gaze and
proximity of many others, acquires not just a collection of beliefs but a
structured landscape of salience: some things matter a great deal, others barely
register.
The cultural-evolution tradition, crystallised in \citet{henrichGilWhite2001Prestige},
has spent the past quarter-century documenting the population-level mechanisms
that produce and maintain these landscapes.
This section surveys that account in enough detail to reveal where a formal
bridge to predictive processing is both possible and necessary.

\subsection*{Prestige as a coordination device for collective attention}

The central conceptual move in \citet{henrichGilWhite2001Prestige} is the
distinction between \emph{prestige} and \emph{dominance} as two structurally
different routes to social influence.
Dominance hierarchies, which appear across social primates, are sustained by
intimidation: individuals defer to dominant animals to avoid costs.
Prestige, by contrast, is \emph{freely conferred deference}.
Copiers voluntarily attend to, sit near, gift, and imitate high-prestige models
because doing so raises their own fitness by providing access to superior
information or skill.
The evolutionary logic is straightforward: when acquiring complex cumulative
culture requires accurately identifying good models, individuals who can cheaply
infer model quality by reading others' deference choices will outcompete those
who must independently assess each candidate.
The community's existing attention distribution therefore functions as a
\emph{second-order prior} over whom it is worth attending to — a public signal
that learners consult before committing their own attentional resources.

\citet{henrichChudekBoyd2015BigMan} formalise this logic in a
culture--gene coevolutionary model of ``big man'' leadership.
In small-scale egalitarian societies, the big man acquires followers not through
coercion but through accumulated demonstrations of generosity, oratory, and
successful co-ordination.
Critically, followers copy big-man behaviours (including prosocial ones), which
in turn creates selection pressure for big men who are genuinely prosocial —
because their only currency is the continued attention and copying their conduct
commands.
The model shows that this prestige channel produces phenotype correlation between
leaders and followers and generates prosocial norms as stable downstream
outcomes: the attention asymmetry is not merely a readout of pre-existing value
distributions but actively reshapes them.
Egalitarian levelling mechanisms — gossip, mockery, ostracism, and in extremis
execution — function not as alternatives to prestige dynamics but as its
complement: they \emph{withdraw} community attention from would-be dominants,
enforcing the condition under which prestige (rather than coercion) remains the
operative currency.
Flat societies are densely attention-managed, not attention-neutral.

The clearest experimental demonstration of how this second-order attention
mechanism works at the level of individual learning is \citet{chudek2012BystanderGaze}.
In a series of studies with 3--4-year-old children, a purely \emph{attentional}
cue — ten seconds of differential bystander gaze toward one of two adult
demonstrators, with no other information about their competence or history — more
than doubled the probability that children subsequently copied that demonstrator
on a novel task.
Crucially, the effect \emph{generalised across domains}: prestige established
through others' gaze in one context (e.g., artifact use) transferred to
unrelated copying choices (e.g., food preferences).
This domain-generality is theoretically decisive.
A content-specific prior — ``this person knows about artifacts'' — would not
transfer to food choices; the prestige cue would stay in the domain where it was
established.
What does transfer is a \emph{precision signal}: ``this person's behaviour is
worth attending to, period.''
In predictive-processing terms (though the authors do not use this vocabulary),
the bystander gaze cue behaves like a generalised precision parameter rather than
a domain-restricted content prior.
It raises the weight accorded to the high-prestige model's actions across the
entire space of observable behaviours, without specifying what those actions will
contain.
Section~\ref{sec:formal} formalises this observation precisely.

\subsection*{CREDs and the costliness logic}

Once language evolved to allow the cheap transmission of beliefs, cultural
learners faced a new inference problem: which professed values does a model
\emph{actually hold}, and which are mere verbal performance?
\citet{henrich2009CREDs} argues that natural selection favoured learners who
weighted models' expressed beliefs by the costliness of their observable
commitment to those beliefs — what he calls Credibility-Enhancing Displays
(CREDs).
Fasting, painful initiations, celibacy, martyrdom, and public sacrifice are
actions whose costs make sense only if the agent genuinely holds the associated
belief or value; they are costly signals that cannot be cheaply faked.
Learners who copied the values of agents who performed CREDs systematically
acquired the values of genuinely committed models, rather than those of
convenient verbal performers.

The fit with precision-weighted inference is near-direct, and Section~\ref{sec:formal}
develops it explicitly.
Informally: without CREDs, a model's verbal claim about a value provides
low-precision evidence — it could reflect genuine commitment or could be
strategic posturing.
A CRED raises the precision of the channel from observed-other-behaviour to
one's own value-prior.
High-CRED demonstrations are high-confidence messages about what the transmitter
actually values; low-CRED environments produce noisy, uncertain updates.
Empirical work by Lanman and Buhrmester \citep{lanmanBuhrmester2017CREDs} has
confirmed this logic: childhood exposure to parental CREDs predicts adult theism
more robustly than competing accounts based on cognitive style or existential
anxiety, suggesting that what gets installed is not a bare belief but a
\emph{calibrated prior} whose strength tracks the evidential quality of the
exposure.
The CRED mechanism and the prestige mechanism are thus complementary channels in
the same installation system: prestige cues direct \emph{which} models receive
attention; CREDs modulate \emph{how much confidence} to accord to those models'
expressed commitments.

\subsection*{Attentional infrastructure is itself culturally inherited}

A natural objection at this point is that while the \emph{contents} of attention
may vary culturally, the \emph{mechanisms} that structure attention — what kinds
of agents and events are even candidates for imitation — must be
evolutionarily fixed.
\citet{heyes2018CognitiveGadgets} mounts a sustained argument that this
objection is wrong.
Her ``cognitive gadgets'' thesis distinguishes a domain-general genetic
``starter kit'' of social attentional biases (sensitivity to faces, voices, and
social rewards) from the higher-order mechanisms of selective social learning
— including imitation, mind-reading, and norm compliance — which are themselves
culturally assembled during development.
The decisive piece of evidence for present purposes: there is empirical support
for the claim that humans can \emph{socially learn attentional biases},
specifically by learning to attend more to some agents than others through the
mechanism of following third-party gaze.
That is, the same bystander-gaze cue that Chudek et al.\ use as an experimental
manipulation of prestige is also the mechanism by which children acquire durable
attentional dispositions — what to find exemplary, whom to watch, which features
of the social environment carry signal.
The salience landscape is not merely populated through cultural transmission; its
\emph{topography} — which gradients channel attention up rather than down — is
itself inherited.

\citet{heyes2024NormPsychology} extends this framework to norm psychology
specifically.
Rather than treating norm-following as a genetically fixed motivational module
(``normative competence as instinct''), she argues that normative competence is a
culturally assembled cognitive gadget built atop domain-general attentional and
motivational capacities.
Crucially, what counts as a norm — which regularities in others' behaviour are
picked up as prescriptive rather than merely descriptive — is not given in
advance but is itself acquired through the same prestige-and-gaze mechanisms
that structure learning more generally.
The implications for any account of value transmission are significant: there is
no innate, domain-general ``norm receptor'' that picks out which social patterns
to treat as obligatory.
Instead, the very categories through which norms are perceived and reproduced
are part of what gets transmitted.
This recursive structure — attentional dispositions shaping what gets copied,
which in turn shapes attentional dispositions in the next generation — is
precisely the loop this paper seeks to formalise.

\subsection*{The ceiling of the cultural-evolution account}

The literature surveyed in this section provides a rich \emph{empirical} story
about value transmission: which individuals attract community attention, under
what conditions, through what cues, with what consequences for the behavioural
repertoires of those who copy them.
The account is anchored in experimental data \citep{chudek2012BystanderGaze},
population-genetic formalisms \citep{henrichChudekBoyd2015BigMan,egoziRam2024PrestigeFormal},
and developmental evidence \citep{heyes2018CognitiveGadgets,heyes2024NormPsychology}.
What this literature does \emph{not} provide is a vocabulary that connects these
dynamics to the computational architecture of the individual learner: why does
bystander gaze function as a domain-general precision signal rather than a
domain-specific content cue?
What is happening, formally, when a CRED raises the confidence with which a
value-prior is adopted?
How do community-level patterns of attention crystallise into the stable
configurations we call norms?
These questions ask for a level of mechanistic description that cultural
evolution's largely informal vocabulary does not supply.
The next section takes up precisely this gap, mapping prestige- and
CRED-mediated attentional dynamics onto the formal language of
precision-weighted predictive processing, and showing how the loop that closes
from community priors through attention and copying back to community priors can
be given a rigorous active-inference rendering.
